% \section{Restrições holonômicas}
% \label{sec:restricoes_holonomicas}

% Devido a vínculos geométricos, nem sempre as coordenadas generalizadas $\vec q$ podem variar livremente. Exemplos incluem:
% (i) um quaternion que deve permanecer unitário, ${}^T q_b q_b=1$;
% (ii) o comprimento fixo de um elo rígido;
% (iii) contato mantido, sem penetração.
% Esses vínculos são chamados de restrições holonômicas e dependem apenas de $\vec q$.

% Para que essas condições sejam satisfeitas em conjunto com a dinâmica, utilizam-se os multiplicadores de Lagrange $(\vec \lambda)$, que representam forças ou torques adicionais introduzidos apenas para garantir o vínculo.
% Essas forças são internas ao modelo matemático e não alteram a física real, apenas garantem o cumprimento das restrições.
% Pode-se modelar o conjunto de $m$ vínculos por:
% \begin{equation}
% \Phi(\vec q,t)=
% \vec 0\in\mathbb {R}^{m},
% \label{eq:PhiJ1}
% \end{equation}

% sendo que:
% \begin{equation}
% \mathbf{J}(\vec q,t)\;\triangleq\;\frac{\partial \Phi(\vec q,t)}{\partial \vec q}\in\mathbb {R}^{m\times n},
% \label{eq:PhiJ2}
% \end{equation}
% Assume-se independência funcional no domínio de operação, isto é, $\operatorname{rank}\mathbf{J}(\vec q,t)=m$.

% Assim, as equações de movimento deixam de ser puramente diferenciais e passam a constituir um sistema diferencial-algébrico (DAE).
% Portanto, a dinâmica com multiplicadores de Lagrange é dada por:
% \begin{equation}
% \mathbf M(\vec q)\,\ddot{\vec q}
% +\mathbf C(\vec q,\dot{\vec q})\,\dot{\vec q}
% +\vec g(\vec q)=
% \vec\tau+\mathbf{J}^{\top}(\vec q,t)\,\vec\lambda,
% \label{eq:dinamica_mult_lagrange}
% \end{equation}
% onde:
% \begin{itemize}
%     \item $\mathbf{M}(\vec{q})$: matriz de inércia;
%     \item $\mathbf{C}(\vec{q},\dot{\vec{q}})$: termos de Coriolis e centrífugos;
%     \item $\mathbf{g}(\vec{q})$: vetor de forças generalizadas de gravidade;
%     \item $\vec{\tau}$: torques/forças generalizadas aplicadas pelos atuadores;
%     \item $\mathbf{J}(\vec q,t)$: matriz Jacobiana das restrições $\Phi(\vec q)$ em relação a $\vec q$;
%     \item $\vec{\lambda}$: multiplicadores de Lagrange, associados às forças ou torques de vínculo.
% \end{itemize}
% e esses termos, acima mencionados, são acoplados ao vínculo algébrico $\Phi(\vec q,t)=\vec 0$, que permite que as restrições holonômicas sejam satisfeitas ao longo da evolução do sistema.

% \subsection{Níveis diferencial e de aceleração}
% Ao derivar
% $\Phi(\vec q,t)=\vec 0$
% no tempo, obtêm-se as condições de velocidade:
% \begin{equation}
% \mathbf{J}(\vec q,t)\,\dot{\vec q}
% +\frac{\partial \Phi(\vec q,t)}{\partial t}
% =\vec 0,
% \label{eq:vellevel}
% \end{equation}

% e ao derivar pela segunda vez, obtêm-se as condições de aceleração:
% \begin{equation}
% \mathbf{J}(\vec q,t)\,\ddot{\vec q}
% +\dot{\mathbf{J}}(\vec q,\dot{\vec q},t)\,\dot{\vec q}+\frac{\partial^2 \Phi(\vec q,t)}{\partial t^2}
% =\vec 0.
% \label{eq:acclevel}
% \end{equation}

% Combinando \eqref{eq:dinamica_mult_lagrange} e
% \eqref{eq:acclevel},
% obtém-se o sistema linear em blocos para
% $\ddot{\vec q}$ e $\vec\lambda$:
% \begin{equation}
% \begin{bmatrix}
% \mathbf M(\vec q)   & -\mathbf{J}^{\top}(\vec q,t)\\
% \mathbf{J}(\vec q,t) & \mathbf 0
% \end{bmatrix}
% \begin{bmatrix}
% \ddot{\vec q}\\
% \vec\lambda
% \end{bmatrix}
% =
% \begin{bmatrix}
% \vec\tau
% -\big(\mathbf C(\vec q,\dot{\vec q})\,\dot{\vec q}
% +\vec g(\vec q)\big)
% -\dot{\mathbf{J}}(\vec q,\dot{\vec q},t)\,\dot{\vec q}-\dfrac{\partial^2 \Phi(\vec q,t)}{\partial t^2}
% \end{bmatrix}.
% \label{eq:sist_linear_blocos}
% \end{equation}
% A equação \eqref{eq:sist_linear_blocos} trata de
% $n{+}m$ equações para $n{+}m$ incógnitas,
% onde:
% \begin{equation}
% \ddot{\vec q}\in\mathbb {R}^{n},
% \qquad
% \vec\lambda\in\mathbb {R}^{m}
% \end{equation}


% As condições iniciais requerem:
% \begin{equation}
% 	\label{eq:ci_sist_linear_blocos1}
% \Phi(\vec q(0),t_0)=\vec 0 
% \end{equation}

% e

% \begin{equation}
% 	\label{eq:ci_sist_linear_blocos2}
% \mathbf{J}(\vec q(0),t_0)\dot{\vec q}(0)
% +\frac{\partial\Phi}{\partial t}\big|_{(\vec q(0),t_0)}=\vec 0.
% \end{equation}


%==================================================
%----- Baumgarte
%==================================================
% \section{Estabilização de Baumgarte}
% Na simulação numérica, mesmo impondo que $\Phi(\vec q) = 0$, os erros de integração podem fazer com que os vínculos se degradem. 
% Por exemplo, em vez de manter
% ${}^T q_b q_b = 1$,
% o quaternion pode tender para ${}^T q_b q_b = 1,0001$.
% A técnica de estabilização de Baumgarte consiste em adicionar uma realimentação proporcional--derivativa no erro do vínculo para permitir que o erro $\Phi$ obedeça a uma equação diferencial estável do tipo:

% \begin{equation}
% \ddot{\Phi}
% +2\,\alpha\,\dot{\Phi}
% +\beta^2\,\Phi
% =\vec 0,
% \label{eq:Baumgarte}
% \end{equation}

% Onde os parâmetros $\alpha$ e $\beta$ são parâmetros do projeto (ganhos de estabilização) escolhidos de forma a amortecer e corrigir desvios do vínculo ao longo do tempo.
% Além disso, quando $\Phi = 0$, pode-se dizer que o vínculo foi satisfeito;
% por sua vez, os termos
% $\dot \Phi$ e $\ddot \Phi$
% são as derivadas (primeira e segunda) do erro.
% Em termos de aceleração, substituindo \eqref{eq:acclevel} por:
% \begin{equation}
% \mathbf{J}(\vec q,t)\,\dot{\vec q}
% =
% -\dot{\mathbf{J}}(\vec q,\dot{\vec q},t)\,\dot{\vec q}
% -\frac{\partial^2 \Phi}{\partial t^2}
% -2\alpha\Big(\mathbf{J}\,\dot{\vec q}
% +\frac{\partial \Phi}{\partial t}\Big)
% -\beta^2\,\Phi.
% \label{eq:acclevel_baumgarte}
% \end{equation}

% A equação \eqref{eq:acclevel_baumgarte} funciona como um amortecedor; se o vínculo se desvia, a equação de Baumgarte conduz a trajetória de volta ao vínculo.
% Uma escolha típica dos parâmetros é
% $\alpha=\zeta\,\omega_n$,
% $\beta=\omega_n$, com
% $\zeta\in[0{,}7;\,1{,}0]$ e
% $\omega_n>0$.
% Sendo que $\zeta$ é a razão de amortecimento;
% o termo $\omega_n$ representa a frequência natural desejada para a dinâmica artificial de estabilização.

% \subsection{Equivalência com controle PD sobre o vínculo.}
% Definindo o erro de vínculo como:
% \begin{equation}
% e(t)\triangleq \Phi(\vec q,t),
% \end{equation}
% a dinâmica de Baumgarte apresentada em \eqref{eq:Baumgarte}
% é a de um sistema de 2ª ordem, e é equivalente a aplicar uma realimentação proporcional--derivativa (PD) no erro do vínculo:
% \begin{equation}
% u_{\Phi}=
% -K_d\,\dot{\Phi}(\vec q,\dot{\vec q},t)
% -K_p\,\Phi(\vec q,t),
% \qquad
% K_d = 2\alpha,\;\; K_p = \beta^2,
% \label{eq:PD_phi}
% \end{equation}
% onde:
% \begin{equation}
% \dot{\Phi}(\vec q, \dot{\vec q}, t) = \mathbf{J}(\vec q,t)\,\dot{\vec q} + \frac{\partial \Phi}{\partial t}.
% \end{equation}

% No nível de aceleração, \eqref{eq:PD_phi} pode ser inserido em:
% \begin{equation}
% \mathbf{J}(\vec q,t)\,\ddot{\vec q}
% = -\,\dot{\mathbf{J}}(\vec q,\dot{\vec q},t)\,\dot{\vec q}
% - \frac{\partial^2 \Phi}{\partial t^2}
% + u_{\Phi},
% \end{equation}
% que é exatamente \eqref{eq:acclevel_baumgarte}, sendo expresso de forma matemática:
% \begin{equation}
% u_{\Phi}=
% -2\alpha\big(\mathbf{J}\dot{\vec q}
% +\frac{\partial \Phi}{\partial t}\big)
% -\beta^2 \Phi.
% \end{equation}

% Para o caso vetorial ($m>1$), usa-se $K_d=2\alpha\,\mathbf I_m$ e $K_p=\beta^2 \mathbf I_m$,
% ou ainda, as matrizes diagonais positivas definidas de ponderação.
% Na parametrização $(\zeta,\omega_n)$,
% com $\alpha=\zeta\omega_n$ e $\beta=\omega_n$, logo:
% \begin{equation}
% K_d = 2 \zeta \omega_n,
% \qquad
% K_p = \omega_n^2
% \end{equation}

% \subsubsection{Critérios de escolha dos ganhos de Baumgarte}
% A seleção dos ganhos deve respeitar certos parâmetros, como:
% \begin{enumerate}[label=(\roman*)]
% \item Amortecimento: tomar $\zeta \in [0{,}7,\,1{,}0]$ para evitar sobressinal de $e(t)=\Phi$ e convergência lenta;
% \item Tempo desejado de acomodação (critério de $2\%$): $t_s \approx \dfrac{4}{\zeta\,\omega_n}$.
% Dado $t_s$ de projeto, escolher $\omega_n = 4/(\zeta\,t_s)$;
% \item Escala e unidades: se $\Phi$ tem unidades físicas (m, rad, etc.), normalizar
% $\tilde{\Phi}=\Phi/\Phi_{\mathrm{ref}}$ para que $K_p=\beta^2$ e $K_d=2\alpha$ não fiquem desbalanceados;
% \item Compatibilidade numérica com o passo $\Delta t$: garantir resolução do modo artificial impondo
% $\omega_n\,\Delta t \lesssim 0{,}3$ (pelo menos 10–20 passos por período $T_\Phi=2\pi/\omega_n$);
% \item Rigidez: valores muito altos de $\omega_n$ tornam o DAE mais rígido e forçam $\Delta t$ pequeno; preferir integradores implícitos quando $\omega_n$ precisar ser elevado.
% \end{enumerate}

% \subsubsection{Impacto na rigidez numérica}
% A elevação de $\omega_n$ desloca autovalores do sistema aumentado para a esquerda (modos rápidos), caracterizando rigidez;
% métodos explícitosimpõem $\Delta t \ll 1/\omega_n$; em métodos $A$-estáveis (por exemplo, trapezoidal) a restrição é atenuada, porém passos muito grandes degradam a precisão e podem reintroduzir desvios via erro de truncamento.

%==================================================
%----- Dissipação
%==================================================
% \section{Dissipação e regularização}
% Mesmo com a estabilização de Baumgarte, podem ocorrer inconsistências numéricas, em especial perto de configurações em que a Jacobiana $\mathbf J(\vec q,t)$ fica mal condicionada. Uma estratégia simples é introduzir amortecimento viscoso linear por meio da \emph{função de dissipação de Rayleigh}:
% \begin{equation}
% \mathcal R(\dot{\vec q}) \;\triangleq\; \tfrac{1}{2}\,\dot{\vec q}^{\top}\,\mathbf D\,\dot{\vec q},
% \qquad \mathbf D=\mathbf D^{\top}\succeq \mathbf 0,
% \label{eq:Rayleigh}
% \end{equation}
% da qual seguem os esforços dissipativos generalizados
% \begin{equation}
% \vec Q_{\!d} \;=\; -\frac{\partial \mathcal R}{\partial \dot{\vec q}}
% \;=\; -\,\mathbf D\,\dot{\vec q}.
% \label{eq:Qd}
% \end{equation}

% \subsection{Energia}
% Se $E=T+V$ é a energia mecânica (com $T$ e $V$ inalterados), então
% $\dot E \;=-\dot{\vec q}^{\top}\mathbf D\,\dot{\vec q}\;=\;-\,2\,\mathcal R(\dot{\vec q}) \;\le 0,$
% isto é, o termo dissipativo retira energia do sistema a uma taxa não negativa.

% \subsection{DAE aumentado com dissipação}
% Acrescentando \eqref{eq:Qd} à dinâmica com multiplicadores,
% \begin{equation}
% \mathbf M(\vec q)\,\ddot{\vec q}
% +\mathbf C(\vec q,\dot{\vec q})\,\dot{\vec q}
% +\mathbf D\,\dot{\vec q}
% +\vec g(\vec q)
% =\vec\tau+\mathbf J^{\top}(\vec q,t)\,\vec\lambda,
% \qquad
% \Phi(\vec q,t)=\vec 0.
% \label{eq:dae_diss}
% \end{equation}
% No nível de aceleração, o sistema em blocos torna-se
% \begin{equation}
% \begin{bmatrix}
% \mathbf M(\vec q) & -\mathbf J^{\top}(\vec q,t)\\
% \mathbf J(\vec q,t) & \mathbf 0
% \end{bmatrix}
% \begin{bmatrix}
% \ddot{\vec q}\\ \vec\lambda
% \end{bmatrix}
% =
% \begin{bmatrix}
% \vec\tau-\mathbf C(\vec q,\dot{\vec q})\dot{\vec q}-\vec g(\vec q)-\mathbf D\,\dot{\vec q}
% -\dot{\mathbf J}(\vec q,\dot{\vec q},t)\,\dot{\vec q}-\dfrac{\partial^2 \Phi}{\partial t^2}
% \end{bmatrix}.
% \label{eq:block_diss}
% \end{equation}

% \subsection{Escolhas práticas na regularização}
% É usual adotar
% $\mathbf D=\operatorname{diag}(d_1,\ldots,d_n)\succeq 0$
% com escalonamento compatível com as unidades de $\vec q$ e com o passo $\Delta t$ do integrador; uma regra pragmática é manter $\omega_{\text{diss}}\Delta t \lesssim 0{,}3$ (10–20 passos por período do modo amortecido).
% O termo $-\mathbf D\dot{\vec q}$ age como amortecedor \emph{numérico}, reduzindo oscilações e melhorando o condicionamento do sistema \eqref{eq:block_diss}. Valores excessivos de $\mathbf D$ podem aumentar a rigidez numérica; se necessário, combinar com métodos implícitos $A$-estáveis. Em problemas com $\mathbf J$ próxima de singularidade, $\mathbf D\succeq 0$ atua em complemento à estabilização de Baumgarte (que age sobre $\Phi,\dot\Phi$), mitigando desvios sem alterar a modelagem de $T$ e $V$.

